\documentclass[10pt,a4paper]{amsart}
\usepackage[margin=19mm]{geometry}
\usepackage{amsmath,amssymb,mathtools}
\usepackage[scr=rsfs,cal=euler]{mathalfa}
\usepackage{xfrac}
\usepackage[unicode,hidelinks,bookmarks=false]{hyperref}
\newcommand{\ie}{i.e.\ }
\newcommand{\cf}{cf.\ }
\newcommand{\resp}{respectively\ }
\newcommand{\Subsection}{\subsection}
\providecommand{\nobreakdash}{\nobreak\hbox{-}}
\setlength{\emergencystretch}{2em}
\multlinegap0pt
\usepackage{bm}
\usepackage{bbm}
\usepackage{dsfont}
\usepackage{xspace}
\usepackage{cancel}
\usepackage{mathtools}
\usepackage{amsthm}
\makeatletter
\@ifundefined{prop}{\newtheorem{prop}{Proposition}}{}
\makeatother
\newcommand{\lt}{\left}
\newcommand{\rt}{\right}
\newcommand{\wt}{\widetilde}
\title[Classification of dynamics for a two person model of planned]{Classification of dynamics for a two person model of planned behavior}
\author{Rishi Dadlani, John S. McAlister, Tahra L. Eissa, Nina H. Fefferman}
\date{Source: arXiv:2601.11804v1 (CC BY 4.0)}
\begin{document}
\maketitle

\noindent\emph{Source and licence.} Classification of dynamics for a two person model of planned behavior. Version arXiv:2601.11804v1 (CC BY 4.0), distributed under the Creative Commons Attribution 4.0 International licence, as recorded in the arXiv licence metadata for this exact version and shown on the abstract page https://arxiv.org/abs/2601.11804v1. Full rights notice: licenses/arxiv-2601.11804-CC-BY-4.0.txt.\par\smallskip

\noindent\textit{Editorial scope.} Counted unit: the Prop whose source label is \texttt{Prop:explicitSolution} in the article (source0.tex lines 324--328, source label \texttt{Prop:explicitSolution}), delivered together with the article's own complete original proof of it (source0.tex lines 329--360, one \texttt{proof} environment of 32 lines). No mathematics is written, repaired or re-derived: statement and proof are the article's own text line for line. Required context retained uncounted: none. Every definition, display and label of those lines is retained unchanged. Omitted from this package and not redistributed: 1--323, 361--883. Nothing inside a retained excerpt is omitted or altered: every retained body is delivered exactly as the article's lines are written, so no omission receipt of any kind accompanies this package. Citations: the retained excerpts carry no citation command at all, so no bibliography entry of the article is transcribed and no reference is redistributed. Reference adaptation: none was needed and none was made; every reference inside the delivered unit resolves inside it, so no unresolved reference or citation is produced. Printed slips kept literal: no printed word, symbol, hypothesis or number of the counted statement or of its proof was altered. Layout and class: the article's own class is \texttt{elsarticle} with packages including amssymb, amsmath, amsthm, mathtools, graphicx; the wrapper is the lane's pinned \texttt{amsart} editorial wrapper at 10pt on A4, which restates the theorem-like environments the retained text needs and adds the article's own macro definitions that the retained text needs (\texttt{\textbackslash lt}, \texttt{\textbackslash rt}, \texttt{\textbackslash wt}), with extra wrapper packages (bm, bbm, dsfont, xspace, cancel, mathtools). The delivered numbering is the unit's own. Assets: no retained span contains a figure environment, a graphics-inclusion command, a PDF-page inclusion, a table or a TikZ picture; no image file is copied into this package, the package declares no asset dependency and no image file is redistributed with it. Licence: version arXiv:2601.11804v1 of this article is distributed under the Creative Commons Attribution 4.0 International licence, as recorded in the arXiv licence metadata for this exact version and shown on the abstract page https://arxiv.org/abs/2601.11804v1. A full rights notice is delivered at licenses/arxiv-2601.11804-CC-BY-4.0.txt. Characters: every retained excerpt is pure ASCII, so the delivered engine, XeLaTeX with its default Latin Modern text font, reports no missing character.
\begin{prop}\label{Prop:explicitSolution}
    Suppose $(h_1), (h_3), (h_4)$. Then, we can write $x_2$'s position after $t$ time, given a starting position of $u_0$ at starting time $kT$, as $x_2(t, u_0) = \tanh(At+Be^{-rt} + Ce^{-2rt} + D(u_0))$. This is valid in $[kT, \min((k+1)T, t^*)]$, where $t^* > t'$, $k \geq 1$ and $x_2(t^*) = \tau$.

    $A$ is as in Proposition 2, $B=-\frac{1}{r}[(\sigma_A\alpha_2-\sigma_S\mu_S)\sigma_C+\sigma_S\sigma_C\mu_C]$, $C=-\frac{\sigma_S\sigma_C}{2r}$, and $D(u_0)=\tanh^{-1}(u_0)-B-C$.
\end{prop}

\begin{proof} In the two individual system, immediately after $x_1$ acts, the system is
\begin{equation}
    \begin{split}
        \dot{x}_1&=[\sigma_A\alpha_1+\sigma_S(y_2-\mu_S)][\sigma_C(y_2+\mu_C)](1-x_1^2)\\
        \dot{x}_2&=[\sigma_A\alpha_2+\sigma_S(y_1-\mu_S)][\sigma_C(y_1+\mu_C)](1-x_2^2)\\
        \dot y_1 &= -ry_1\\
        \dot y_2 &= -ry_2
    \end{split}
\end{equation}
And the initial conditions are $x_1(0)=0,\; x_2(0)= u_0,\; y_1(0)=1,\;y_2(0)= \wt y $, where $u_0$ is the position of $x_2$ when $x_1$ acts, and $\wt y$ similarly. Notice that for our constraints, the system is decoupled, so we need only consider $x_2$ and $y_1$. Furthermore $y_1(t)=e^{-rt}$ so we have a single separable ODE. The computation follows through grouping and keeping track of the constants:
\[\dot{x}_2=[C_1+C_2e^{-rt}+C_3][C_4e^{-rt}+C_5](1-x_2^2)\]
where $C_1=\sigma_A\alpha_2, \;C_2 = \sigma_S,\; C_3=-\sigma_S\mu_S,\; C_4=\sigma_C$, and $C_5=\sigma_C\mu_C$. Simplifying:
\[\dot x_2 = [C_6+C_7e^{-rt}+C_8e^{-2rt}](1-x_2^2)\]
where $C_6= A_2 = (\sigma_A\alpha_2-\sigma_S\mu_S)(\sigma_C\mu_C),\; C_7 = (\sigma_A\alpha_2-\sigma_S\mu_S)\sigma_C+ \sigma_S\sigma_C\mu_C,\; C_8=\sigma_S\sigma_C$. For brevity we write $x_2$ just as $x$.
\begin{equation*}
    \begin{split}
        \frac{1}{1-x^2}\frac{dx}{dt}&=C_6+C_7e^{-rt}+C_8e^{-2rt}\\
        \int\frac{1}{1-x^2}\frac{dx}{dt}dt&=\int(C_6+C_7e^{-rt}+C_8e^{-2rt})dt\\
        \int\frac{1}{1-x^2}{dx}&=\big[C_6t-\frac{C_7}{r}e^{-rt}-\frac{C_8}{2r}e^{-2rt}\big]+C_0
    \end{split}
\end{equation*}
where $C_9=\frac{1}{r}[(\sigma_A\alpha_2-\sigma_S\mu_S)\sigma_C+\sigma_S\sigma_C\mu_C]$, $C_{10}=\frac{\sigma_S\sigma_C}{2r}$, and $C_0$ is the constant of integration. Through partial fraction decomposition on the LHS, we obtain the following
\[\frac{1}{2}\ln\lt(\frac{1+x}{1-x}\rt)=\big[C_6t-\frac{C_7}{r}e^{-rt}-\frac{C_8}{2r}e^{-2rt}\big]+C_0\]
As $x$ is bounded in $(-1,1)$, we need not use the absolute values. Using this, we can solve for $C_0(u_0)$ through $x(0)=u_0$, which gives us
\[C_0(u_0)=\frac{1}{2}\ln\lt(\frac{1+u_0}{1-u_0}\rt)+\frac{C_7}{r}+\frac{C_8}{2r}\]
We rearrange the equation again to get that
\[\frac{1+x}{1-x}=\exp\Big(2\big(\underbrace{C_{6}t-\frac{C_{7}}{r}e^{-rt}-\frac{C_{8}}{2r}e^{-2rt}+C_0(u_0)}_{h(t,u_0)}\big)\Big)\]
The left-hand side of the function is injective from $(-1,1)\to(0,\infty)$ so we can solve for the unique solution
\[x=\frac{e^{2h(t,u_0)}-1}{e^{2h(t,u_0)}+1}=\tanh\lt(h(t,u_0)\rt)\]
So the solution is written as follows
\[x(t;u_0) =\tanh(At+Be^{-rt}+Ce^{-2rt}+D(u_0))\]
where $A= A_2 = (\sigma_A\alpha_2-\sigma_S\mu_S)(\sigma_C\mu_C),\; B=-\frac{1}{r}[(\sigma_A\alpha_2-\sigma_S\mu_S)\sigma_C+\sigma_S\sigma_C\mu_C],\; C=-\frac{\sigma_S\sigma_C}{2r},\; D(u_0)=\frac{1}{2}\ln\lt(\frac{1+u_0}{1-u_0}\rt)-B-C=\tanh^{-1}(u_0)-B-C$ \end{proof}

\end{document}
